\section{Experiments}
\label{sec:experiments}

The study has 24 runs: six base runs, 18 ablation runs on DistilBERT and
three runs of the winning configuration on BERT, all on one RTX A6000.

\subsection{Performance}

\begin{table}[t]
\centering
\caption{Base runs with each library's default head. Test metrics in \%.
Precision, recall and F1 are macro averages. Latency and memory come from the
controlled benchmark at batch 1 and the dataset's input length. Bold marks the
better model per dataset.}
\label{tab:results}
\small\setlength{\tabcolsep}{3.6pt}
\tabla{tabla_resultados}
\end{table}

Table~\ref{tab:results} holds the six base runs. BERT wins on SST-2 by 1.04
accuracy points and on Yelp by 0.46. On AG News the sign flips to
DistilBERT by 0.09 points, which is seven test items out of 7{,}600. We read
that as a tie. DistilBERT's 91.28\% on SST-2 matches the 91.3 reported by
\citet{sanh2019distilbert}, a sign that the splits leak nothing and the recipe
is sound.

\begin{figure}[t]
\centering
\includegraphics[width=\linewidth]{curvas_todas.pdf}
\caption{Training loss (logged every 100 steps) and validation loss (every
500 steps) against optimizer iterations for the six base runs.}
\label{fig:curves}
\end{figure}

In Figure~\ref{fig:curves} both models trace almost the same loss path. On
SST-2, training loss keeps falling in the second epoch while validation loss
stalls near 0.17, which is why we select checkpoints on validation F1.

\subsection{Efficiency}
\label{sec:efficiency}

\begin{figure}[t]
\centering
\begin{minipage}[t]{0.49\linewidth}
  \centering
  \includegraphics[width=\linewidth]{burbujas_params_acc.pdf}
  \caption{Parameters against test accuracy. Bubble area is latency at
  batch 1 and the dataset's input length.}
  \label{fig:bubbles}
\end{minipage}\hfill
\begin{minipage}[t]{0.49\linewidth}
  \centering
  \includegraphics[width=\linewidth]{ablation.pdf}
  \caption{Validation F1 of the six DistilBERT configurations. The lilac row
  is the selected one. The right column is the mean over datasets.}
  \label{fig:ablation}
\end{minipage}
\end{figure}

\begin{figure}[t]
\centering
\includegraphics[width=\linewidth]{benchmark_eficiencia.pdf}
\caption{Controlled benchmark (RTX A6000, fp32, one process). Three latency
panels, one per batch size, with the mean speed-up of DistilBERT written in
each. The last panel shows peak GPU memory at batch 1 and batch 32.}
\label{fig:benchmark}
\end{figure}

Figure~\ref{fig:bubbles} puts size, accuracy and latency in one picture. The
DistilBERT column sits 42.5\,M parameters to the left of BERT with bubbles
half as large, and the vertical distance between each pair is the accuracy
price. On the controlled grid (Figure~\ref{fig:benchmark}) DistilBERT is
faster at all nine points, by $1.86\times$ to $2.15\times$. At batch 32 and
64 tokens it processes 2{,}535 samples per second against 1{,}282, and its peak
memory is 62\% of BERT's at batch 1 and 77\% at batch 32 with 256 tokens.
Training the three datasets took 16.4 minutes against 29.3.

\subsection{Ablation study}

\begin{table}[t]
\centering
\caption{Ablation on DistilBERT. Validation F1 in \%. The lilac row is the
selected configuration (highest mean), and bold marks the best value per column.
Time is the sum over the three datasets.}
\label{tab:ablation}
\small\setlength{\tabcolsep}{3.6pt}
\tabla{tabla_ablation}
\end{table}

Table~\ref{tab:ablation} and Figure~\ref{fig:ablation} cover the six
configurations. The five heads with a trainable encoder span 95.03 to 95.16
mean validation F1, and their order is unstable across datasets
(\texttt{h768} is the best on SST-2 and the worst on average). The frozen
encoder drops to 88.87. Against \texttt{h768}, which has the same head, it
loses 9.1 points on SST-2, 5.9 on Yelp and 3.4 on AG News, in exchange for
training 0.6\,M parameters in 6.0 minutes. The linear head wins on the mean,
and it is also the smallest and fastest to train.

\subsection{Best configuration on BERT}

\begin{table}[t]
\centering
\caption{Linear head (\texttt{lineal}) on both models, same recipe. Test
metrics and validation F1 in \%, efficiency as in Table~\ref{tab:results}.
Bold marks the better model per dataset.}
\label{tab:best}
\small\setlength{\tabcolsep}{3.6pt}
\tabla{tabla_mejor_config}
\end{table}

With the linear head (Table~\ref{tab:best}) BERT wins all three datasets,
by 1.95 points on SST-2, 0.43 on Yelp and 0.11 on AG News.
